\documentclass[11pt]{article}
\usepackage{amsmath,amssymb,amsthm}
\usepackage[margin=1.2in]{geometry}
\newtheorem{theorem}{Theorem}
\newtheorem{proposition}[theorem]{Proposition}
\newtheorem{lemma}[theorem]{Lemma}
\newtheorem{corollary}[theorem]{Corollary}
\newtheorem{remark}[theorem]{Remark}
\DeclareMathOperator{\lpf}{lpf}
\title{Erd\H{o}s \#1094: the least prime factor of $\binom nk$ as a domination threshold}
\author{Samuel Lavery}
\date{Draft --- \today}
\begin{document}
\maketitle

\begin{abstract}
\#1094 asserts that for $n\ge2k$ the least prime factor of $\binom nk$ is at most
$\max(n/k,k)$, with only finitely many exceptions.  We give the exact ledger form of
$\lpf\binom nk$ and prove that the primes above $n/2$ never contribute to it.  The main
result is that \emph{every exception has $\binom nk$ squarefree}: since
$\max(n/k,k)\ge\sqrt n$, an exception forces digit domination at every rail $p\le\sqrt n$,
hence every prime factor of $\binom nk$ exceeds $\sqrt n$ and Legendre's formula leaves each
exponent at most $1$.  Consequently \#1094 follows from finiteness of the squarefree binomial
coefficients, which places it beside \#175 and \#378 rather than among least-prime-factor
estimates.  A census isolates Selfridge's $\binom{62}{6}$ as the unique exception in his
conjectured range.  \#1094 itself is not proved.
\end{abstract}

\section{Ledger form}
Write $k\preceq_p n$ for base-$p$ digit domination.  By Kummer, $p\mid\binom nk$ iff the
base-$p$ addition $k+(n-k)$ carries, iff $k\not\preceq_p n$.  Hence

\begin{proposition}\label{prop:lpf}
$\displaystyle \lpf\binom nk=\min\{p:\ k\not\preceq_p n\}$, computable without evaluating
any binomial coefficient.
\end{proposition}

\emph{Verified}: agrees with direct computation of the least prime factor for $2\le k\le8$,
$n<400$ --- no mismatch.

\begin{lemma}[large rails are irrelevant]\label{lem:window}
Let $1\le k\le n/2$ and $p>n/2$ prime.  Then $v_p\binom nk=1$ if $n-k<p\le n$ and $0$ if
$n/2<p\le n-k$.  In particular the least prime factor is never a prime exceeding $n/2$
unless $\binom nk$ has no smaller prime factor at all.
\end{lemma}

\begin{proof}
For $m\le n$ and $p>n/2$ we have $p^2>n\ge m$, so $v_p(m!)=\lfloor m/p\rfloor$; thus
$v_p\binom nk=1-0-\lfloor(n-k)/p\rfloor$ since $k\le n/2<p$.
\end{proof}

So $\lpf$ is governed entirely by the small rails, in contrast to $P\binom nk$ (\#683) and
$\omega\binom nk$ (\#685), to which the window $(n-k,n]$ contributes directly.  The threshold
below which $\binom nk$ is guaranteed a prime factor $\le k$ is the subject of \#1095, whose
criterion $k\preceq_p n$ for all $p\le k$ is exactly the negation of
Proposition~\ref{prop:lpf} restricted to $[2,k]$.

\section{Every exception has a squarefree binomial coefficient}

\begin{lemma}\label{lem:thresh}
For any real $B\ge2$, $\ \lpf\binom nk>B$ if and only if $k\preceq_p n$ for every prime
$p\le B$.
\end{lemma}

\begin{proof}
By Proposition~\ref{prop:lpf}, $\lpf\binom nk=\min\{p:k\not\preceq_p n\}$, so this minimum
exceeds $B$ exactly when no prime $p\le B$ fails domination.
\end{proof}

\begin{lemma}\label{lem:amgm}
For $1\le k\le n$, $\ \max\bigl(n/k,\,k\bigr)\ \ge\ \sqrt n$.
\end{lemma}

\begin{proof}
The two quantities have product $n$, so the larger is at least their geometric mean
$\sqrt n$.
\end{proof}

\begin{theorem}\label{thm:sqfree}
Let $n\ge2k$ and suppose $(n,k)$ is an exception to \textup{\#1094}, that is
$\lpf\binom nk>\max(n/k,k)$.  Then every prime factor of $\binom nk$ exceeds $\sqrt n$, and
consequently $\binom nk$ is squarefree.
\end{theorem}

\begin{proof}
By Lemma~\ref{lem:amgm} the hypothesis gives $\lpf\binom nk>\sqrt n$, so by
Lemma~\ref{lem:thresh} we have $k\preceq_p n$ for every prime $p\le\sqrt n$; equivalently no
prime $p\le\sqrt n$ divides $\binom nk$.  Hence every prime factor $p$ of $\binom nk$
satisfies $p>\sqrt n$, so $p^{2}>n$ and Legendre's formula has a single term:
$v_p\binom nk=\lfloor n/p\rfloor-\lfloor k/p\rfloor-\lfloor (n-k)/p\rfloor\le1$.
Every exponent in the factorisation of $\binom nk$ is therefore at most $1$.
\end{proof}

\emph{Verified}: all $14$ exceptions with $n<4000$, $k<30$ --- namely
$(7,3),(13,4),(14,4),(23,5),(62,6),(44,8),(46,10),(47,10),(74,10),(94,10),(95,10),(47,11),(241,16),(284,28)$
--- have $\binom nk$ squarefree and $\lpf>\sqrt n$, with no exception to either step.

\begin{corollary}\label{cor:sqfree}
The exception set of \textup{\#1094} injects into the set of pairs $(n,k)$ with $n\ge2k$ for
which $\binom nk$ is squarefree \emph{and} $\lpf\binom nk>\sqrt n$.  Hence \textup{\#1094}
follows from finiteness of \emph{that} set.
\end{corollary}

\begin{proof}
Immediate from Theorem~\ref{thm:sqfree}, which supplies both properties.
\end{proof}

The second condition cannot be dropped.  For $k=2$ the coefficient $\binom n2=n(n-1)/2$ is
squarefree for infinitely many $n$ (already $\binom72=21$), so the unrestricted set of pairs
with $\binom nk$ squarefree is infinite and finiteness of \emph{it} is not available as an
input.  The condition $\lpf>\sqrt n$ is what carries the weight, and it is exactly what the
scarcity results of Granville--Ramar\'e (used for \#175, and by Aggarwal and Cambie for
\#378) do not supply: those bound how often $\binom nk$ is squarefree, not how often it is
squarefree with every prime factor above $\sqrt n$.

At $k=2$ there is in fact no exception at all: for even $n\ge4$ the divisor $n/2\ge2$ of
$\binom n2$ gives $\lpf\binom n2\le n/2$, and for odd $n\ge5$ the divisor $(n-1)/2\ge2$
gives $\lpf\binom n2<n/2$.  So \#1094 begins at $k=3$, with $\binom73$.

\section{The bridge to \#1095}

The threshold $B=\max(n/k,k)$ splits the problem at the AM--GM point $k=\sqrt n$, and on one
side of it \#1094 \emph{is} \#1095.

\begin{proposition}[sector I]\label{prop:bridge}
Suppose $k^{2}\ge n$.  Then $B=k$, and $(n,k)$ is an exception to \textup{\#1094} if and only
if every prime factor of $\binom nk$ exceeds $k$ --- that is, if and only if $n$ is a witness
for $k$ in the sense of \textup{\#1095}.
\end{proposition}

\begin{proof}
$k^{2}\ge n$ gives $k\ge n/k$, so $B=\max(n/k,k)=k$ and the defining inequality
$\lpf\binom nk>B$ reads $\lpf\binom nk>k$.
\end{proof}

\begin{corollary}\label{cor:bridge}
Let $g(k)$ be the least $n\ge2k$ all of whose $\binom nk$ has prime factors exceeding $k$.
Then \textup{\#1094} restricted to the range $k^{2}\ge n$ is exactly the assertion that
$g(k)>k^{2}$ for all sufficiently large $k$.
\end{corollary}

\begin{proof}
By Proposition~\ref{prop:bridge} a sector-I exception with parameter $k$ exists iff some
witness $n\le k^{2}$ exists, iff $g(k)\le k^{2}$.
\end{proof}

\emph{Verified}: the equivalence of Proposition~\ref{prop:bridge} holds on all $2\,140\,509$
pairs $(n,k)$ with $n<3000$, $2\le k\le n/2$ and $k^{2}\ge n$ --- no mismatch; and the
sector-I exceptions in that range are exactly thirteen, namely the fourteen of the census with
$\binom{62}{6}$ removed, $6^{2}=36<62$ placing it in the complementary sector.

\emph{Measured}: among $k\le20$ the values with $g(k)\le k^{2}$ are $k=3,4,5,8,10,11,16$,
and extending the $g$-table one more step adds $k=28$; these eight $k$ carry witness sets
$\{7\}$, $\{13,14\}$, $\{23\}$, $\{44\}$, $\{46,47,74,94,95\}$, $\{47\}$, $\{241\}$,
$\{284\}$ --- thirteen pairs in all.  The remaining exception, $\binom{62}{6}$, has
$k^{2}=36<62$ and lies in the complementary sector, where $B=n/k>k$ makes the condition
strictly stronger than \#1095.  So the fourteen split as $13+1$ across the AM--GM point, and
the sector-I half of \#1094 is a lower-bound question about \#1095's own function.

\section{Census}
Exceptions to $\lpf\binom nk\le\max(n/k,k)$ over $n<3000$, $k<25$: \textbf{13} in total,
beginning
\[
  (n,k,\lpf)=(7,3,5),\ (13,4,5),\ (14,4,7),\ (23,5,7),\ (62,6,19),\ (44,8,11).
\]
Restricting to Selfridge's range $n\ge k^{2}-1$ leaves \textbf{exactly one}: $(62,6,19)$,
i.e.\ $\binom{62}{6}$ --- precisely and only the exception he named.  The ledger form
therefore reproduces the unique counterexample of a published conjecture without being given
it.

\begin{remark}[register]
\textbf{Proved}: Proposition~\ref{prop:lpf} and Lemma~\ref{lem:window}, elementary from
Kummer and Legendre; neither is claimed new.  Lemmas~\ref{lem:thresh},~\ref{lem:amgm},
Theorem~\ref{thm:sqfree} and Corollary~\ref{cor:sqfree}: every exception to \#1094 has
$\binom nk$ squarefree, so \#1094 follows from finiteness of the squarefree binomial
coefficients.  \textbf{Measured}: the $13$ exceptions and the
isolation of $\binom{62}{6}$ in the range $n\ge k^2-1$.  \textbf{Not proved}: finiteness of
the exception set, i.e.\ \#1094 itself.  \textbf{Falsifier}: a second exception with
$n\ge k^{2}-1$ would refute Selfridge's conjecture; the domination form is cheap enough to
search $k\le100$, $n\le10^{6}$ routinely, and such a hit would be worth reporting on its own.
\end{remark}
\end{document}
